\chapter{Synthetic Data Engineering}
\label{ch:synthetic}

In the absence of real transaction data --- a common constraint in banking anomaly detection projects --- the quality of the synthetic dataset determines the ceiling of model performance. This chapter covers the source extraction methodology, archetype design, and the generator that produces labeled transactions.

\section{Source Extraction Methodology}

The generator requires two classes of calibration data: parameters describing the \textbf{normal population} (how legitimate clients transact) and patterns describing the \textbf{flagged population} (what suspicious activity looks like). No single source covers both, so a multi-source strategy was designed.

A structured extraction template scores each source across seven trust dimensions --- recency, authority, peer review, information volume, geographic relevance, data type, and reproducibility --- as Strong, Medium, or Weak. A nine-dimension coverage matrix then evaluates each source's contribution: amount distribution, operation mix, frequency, inter-event time, temporal patterns, spatial/branch patterns, channel mix, counterparty known-rate, and \emph{operation sequence patterns}. The last dimension was added specifically for the LSTM, which requires cross-operation temporal ordering.

\section{Source Extractions}

\textbf{S1 --- PaySim.} A synthetic mobile-money simulator based on real transaction logs from an undisclosed African provider. Trust profile: 2 Strong (peer review, reproducibility), 4 Medium, 1 Weak (geographic relevance --- mobile money in sub-Saharan Africa differs structurally from Tunisian branch banking). Its primary role is as a \emph{beneficiary graph structure} source: the agent-to-agent transfer network is the best available model for counterparty relationship patterns. Secondary roles are amount distribution shape and day-of-month seasonality. Operation mapping is direct for three of four operations (\code{CASH\_OUT} $\leftrightarrow$ retrait, \code{CASH\_IN} $\leftrightarrow$ versement, \code{TRANSFER} $\leftrightarrow$ virement); cheque has no PaySim equivalent. The PaySim anomaly base rate of $\sim$0.13\% was retained as the initial injection-rate anchor.

\textbf{S2 --- Amen Bank 2023 Annual Report.} Provides \emph{normal population} calibration. The document has an internal two-tier trust structure: audited financial statements (signed by BDO Tunisie and La G\'en\'erale d'Audit et Conseil) versus unaudited management narrative --- a distinction that guided which figures were treated as hard constraints. Key extractions: the client deposit composition (61\% individuals / 33\% private business / 6\% institutional) calibrating archetype proportions; 150 branches across 9 zones setting the branch pool; 1{,}139 employees with 60.5\% network-assigned, which combined with the branch count yields the approximation of 2 tellers per branch; and confirmation of the existing AI-AML system, validating the project's positioning as complementary.

\textbf{S3 --- CTAF 2021 Annual Report.} Provides \emph{flagged population} calibration. Four typological case studies were extracted with a structured 7-field template: esp\`eces-heavy structuring below the 10{,}000~DT threshold dispersed across branches (informing the near-threshold ratio feature and the smurfing scenario); a virement chain to unknown beneficiaries with round amounts (informing first-time-beneficiary and round-amount rules); a politically exposed person used as a fund conduit (validating the decision to treat PPE as a structural flag layered on archetypes rather than a separate archetype); and third-party cheques deposited in rapid succession (informing cumulative amount tracking). Six quantitative targets were also extracted, including instrument co-occurrence rates and the concentration of BBE operations in the 5{,}000--20{,}000~DT band.

\textbf{Known gap --- BCT circulars.} Circulars 2018-09, 2017-08, and 2021-05 would have provided an exhaustive threshold enumeration for internal LAB/FT controls, but are restricted documents. The gap is documented rather than papered over: the Decision Service's rules derive from CTAF case studies and Law n\textsuperscript{o}41-2024, not a complete circular inventory.

\section{Client Archetype Design}

Five archetypes were designed across six behavioral dimensions, calibrated from the S2 deposit composition.

\begin{table}[htbp]
\centering
\small
\caption{Client archetypes}
\begin{tabular}{@{}lcccc@{}}
\toprule
\textbf{Archetype} & \textbf{Population} & \textbf{Dominant op.} & \textbf{Monthly vol.} & \textbf{Amount range (DT)} \\
\midrule
Salaried & 49\% & Retrait (50\%) & $\sim$5.7 & 200--500 \\
Small Business & 33\% & Versement (55\%) & $\sim$37.5 & 500--2{,}000 \\
Student & 5\% & Retrait (95\%) & $\sim$7.2 & 40--80 \\
Retiree & 7\% & Retrait / Virement & $\sim$3.2 & 200--500 \\
Big Business & 6\% & Virement (78\%) & $\sim$63 & 2{,}000--20{,}000 \\
\bottomrule
\end{tabular}
\end{table}

Proportions derive from the 61/33/6 split, with the individual segment subdivided approximately 80\% salaried, 12\% retiree, 8\% student. The six dimensions per archetype are operation mix, transaction frequency (mean and std per operation per month), timing (preferred day of month), amount ranges, counterparty loyalty, and branch loyalty. All parameters are externalized in a YAML configuration file, enabling recalibration without code changes.

A regulatory discovery during calibration --- Law n\textsuperscript{o}41-2024 capping cheques at 30{,}000~DT, followed by a 24.9\% collapse in cheque volume --- required revising the big-business archetype: \emph{remise de ch\`eques} frequency dropped from $\sim$12--15 per month to 2--4, with \emph{virement} absorbing the difference. A new anomaly scenario, cheque structuring just under the ceiling, was added.

\section{Generator Architecture}

The generator follows a four-class, two-phase design separating the behavioral model from the anomaly injection mechanism:

\begin{enumerate}
  \item \textbf{Archetype} (layer 1) --- loads behavioral parameters from YAML, defining population-level priors.
  \item \textbf{Client} (layer 2) --- samples a fixed personality from the archetype prior via Gaussian sampling with clipping. Each client's parameters are drawn once at creation and remain fixed for life; operation mix probabilities are sampled then normalized to sum to 1.
  \item \textbf{Generator} (orchestrator) --- \emph{Phase 1 (planning)} creates the population, estimates total event volume, and assigns anomaly scenarios to selected clients with specific time windows and difficulty tiers. \emph{Phase 2 (generation)} ticks through each simulated day, determines which clients transact (bell curve around their preferred day), and emits normal or anomalous events depending on whether the client is inside an active anomaly window.
  \item \textbf{Event} (inline) --- per-event noise, payload construction, branch selection respecting branch loyalty, and employee assignment from the visited branch's pool.
\end{enumerate}

The two-phase approach is essential for controlled injection: hard-tier scenarios like smurfing span 21 days and require sufficient normal history \emph{before} the window for the LSTM to have learned the baseline. Planning windows before generation guarantees this temporal constraint.

\subsection{Branch-Linked Employee Assignment}

A data quality requirement worth highlighting: each employee belongs to exactly one branch (2 per branch, 300 total). The initial implementation drew employees from a global pool, allowing the same teller to appear at multiple branches --- physically impossible in branch banking. This contaminated the \code{repeated\_client\_employee} scenario, which requires a specific client to visit a specific teller at a specific branch repeatedly and is meaningless if the teller can appear anywhere. After fixing the assignment, the LSTM's AUC improved from 0.9430 to 0.9866, because branch-linked employees produce cleaner sequential patterns.

\section{Anomaly Scenario Design}

Nine scenarios span three difficulty tiers, where difficulty reflects how many features an anomaly modifies and over what horizon.

\begin{table}[htbp]
\centering
\small
\caption{Anomaly scenarios by difficulty tier}
\begin{tabular}{@{}llp{6.6cm}l@{}}
\toprule
\textbf{Tier} & \textbf{Scenario} & \textbf{Description} & \textbf{Actor} \\
\midrule
\multirow{3}{*}{Easy (1d)}
 & Amount spike & Single transaction 5--10$\times$ client mean & Client \\
 & Round amount & Large round amount (5{,}000--20{,}000~DT) & Client \\
 & Duplicate virement & Matching virement within short window & Client \\
\midrule
\multirow{3}{*}{Medium (5d)}
 & Frequency burst & Transaction rate $>$ 3$\times$ daily average & Client \\
 & Cumulative threshold & 24h/48h cumulative exceeds norm & Client \\
 & Timing anomaly & Transactions outside branch hours & Client \\
\midrule
\multirow{3}{*}{Hard (21d)}
 & Smurfing & Versements in 8{,}000--9{,}800~DT, multiple branches & Client \\
 & Cheque structuring & Cheques near the 30{,}000~DT ceiling & Client \\
 & Repeated client-employee & Same client--teller pair $\geq$3 in 24h & Cross \\
\bottomrule
\end{tabular}
\end{table}

Every scenario was validated for \emph{feature-space detectability}: the injected anomaly must produce a measurable change in at least one of the 30 enriched features. Scenarios that modified only profile-level aggregates were removed during a mid-project audit, reducing the count from 13 to 9. The original \code{behavioral\_drift} scenario, for example, gradually shifted a client's operation mix over 21 days, but the shift was absorbed by the profile's own 30-day distribution window and produced no detectable per-event signal.

The difficulty distribution follows a pyramid --- 10\% easy, 30\% medium, 60\% hard --- reflecting the observation that most suspicious activity is designed to evade detection.

\section{Data Validation}

The generator produces four CSV files per run: \code{transactions.csv} (243{,}904 rows), \code{clients\_master.csv} and \code{accounts.csv} (10{,}000 each), and \code{employees\_master.csv} (300).

\begin{table}[htbp]
\centering
\small
\caption{Dataset validation}
\begin{tabular}{@{}ll@{}}
\toprule
\textbf{Metric} & \textbf{Value} \\
\midrule
Total events & 243{,}904 \\
Anomaly rate & 0.623\% (target $\sim$0.13\%) \\
Anomaly scenarios & 9 (all represented) \\
Unique clients & 10{,}000 across 5 archetypes \\
Employees & 300 (2 per branch, branch-linked) \\
Cross-branch employee violations & 0 \\
Simulation period & 180 days (January--June 2025) \\
Salaried retrait mean & $\sim$300~DT (consistent with Tunisian salary) \\
\bottomrule
\end{tabular}
\end{table}

The realized anomaly rate exceeds the 0.13\% target because multi-day windows (5 days medium, 21 days hard) produce multiple anomalous events per plan entry rather than one. This was accepted: the rate remains well within the imbalance range where reconstruction- and sequence-based models learn effectively, and AUC-ROC is threshold-independent.

\subsection{Oracle Profile Shortcut}

For training, client personality parameters are exported directly from the generator's in-memory objects to \code{clients\_master.csv}, providing the exact behavioral baseline needed for contrast features without a batch profile computation over simulated history. The shortcut produces the same 30 features as the production pipeline, preserving training--inference parity. In production these baselines are replaced by batch-computed profiles from real history --- a substitution the pipeline handles without code changes.
